% Source: https://artofproblemsolving.com/community/c3h3165992_2018_lmt_fall_guts_round_p13p24__lexington_mathematical_tournament
% Stable identity: aops-topic-3165992-q21
Scale and rotate so $C=(0,0)$ and $A=(\sqrt3,1)$. Since the line $BA$ makes angle 15 degrees above the leftward horizontal, $B=(\sqrt3+\cot15^\circ,0)=(2\sqrt3+2,0)$. Rotating $\overrightarrow{AC}=(-\sqrt3,-1)$ counterclockwise by 90 degrees gives $\overrightarrow{AD}=(1,-\sqrt3)$, hence $D=(\sqrt3+1,1-\sqrt3)$, below $BC$ as required. The direction of $\overrightarrow{DA}=(-1,\sqrt3)$ is 120 degrees. The vector $\overrightarrow{DB}=(\sqrt3+1,\sqrt3-1)$ has slope $2-\sqrt3=\tan15^\circ$ and points rightward. Therefore $\boxed{\angle ADB=120^\circ-15^\circ=105^\circ}$.
\begin{asy}
usepackage("amsmath");
size(7cm);
pair C=(0,0),A=(sqrt(3),1),B=(2*sqrt(3)+2,0),D=(sqrt(3)+1,1-sqrt(3));draw(A--B--C--cycle);draw(A--D--C);draw(D--B,blue);label("$A$",A,N);label("$B$",B,E);label("$C$",C,W);label("$D$",D,S);
\end{asy}
